% Ad Atticum 4.1 (ad-atticum-04-01)
% Source: https://raw.githubusercontent.com/PerseusDL/canonical-latinLit/master/data/phi0474/phi057/phi0474.phi057.perseus-lat1.xml
% Fetched by scripts/fetch_latin.py

% Dateline (Perseus): Scr. Romae med. m. Sept. a. 697 (57) .

\ciceroLetterOpener{CICERO ATTICO salutem}

\ciceroSection{1} cum primum Romam veni fuitque cui recte ad te litteras darem, nihil prius faciendum mihi putavi quam ut tibi absenti de reditu nostro gratularer. cognoram enim, ut vere scribam, te in consiliis mihi dandis nec fortiorem nec prudentiorem quam me ipsum nec etiam pro praeterita mea in te observantia nimium in custodia salutis meae diligentem eundemque te, qui primis temporibus erroris nostri aut potius furoris particeps et falsi timoris socius fuisses, acerbissime discidium nostrum tulisse plurimumque operae, studi, diligentiae, laboris ad conficiendum reditum meum contulisse.

\ciceroSection{2} itaque hoc tibi vere adfirmo, in maxima laetitia et exoptatissima gratulatione unum ad cumulandum gaudium conspectum aut potius complexum mihi tuum defuisse. quem semel nactus si umquam dimisero ac nisi etiam praetermissos fructus tuae suavitatis praeteriti temporis omnis exegero, profecto hac restitutione fortunae me ipse non satis dignum iudicabo.

\ciceroSection{3} nos adhuc, in nostro statu quod difficillime reciperari posse arbitrati sumus, splendorem nostrum illum forensem et in senatu auctoritatem et apud viros bonos gratiam magis quam optaramus consecuti sumus; in re autem familiari, quae quem ad modum fracta, dissipata, direpta sit non ignoras, valde laboramus tuarumque non tam facultatum quas ego nostras esse iudico quam consiliorum ad conligendas et constituendas reliquias nostras indigemus.

\ciceroSection{4} nunc etsi omnia aut scripta esse a tuis arbitror aut etiam nuntiis ac rumore perlata, tamen ea scribam brevi quae te puto potissimum ex meis litteris velle cognoscere. Pr. Nonas Sextilis Dyrrachio sum profectus ipso illo die quo lex est lata de nobis. Brundisium ven i Nonis Sextilibus. ibi mihi Tulliola mea fuit praesto natali suo ipso die qui casu idem natalis erat et Brundisinae coloniae et tuae vicinae Salutis; quae res animadversa a multitudine summa Brundinisorum gratulatione celebrata est. ante diem iii Idus Sextilis cognovi, quom Brundisi essem, litteris Quinti mirifico studio omnium aetatum atque ordinum, incredibili concursu Italiae legem comitiis centuriatis esse perlatam. inde a Brundisinis honestissime ornatus iter ita feci ut undique ad me cum gratulatione legati convenerint.

\ciceroSection{5} ad urbem ita veni ut nemo ullius ordinis homo nomenclatori notus fuerit qui mihi obviam non venerit, praeter eos inimicos quibus id ipsum, se inimicos esse, non liceret aut dissimulare aut negare. cum venissem ad portam Capenam, gradus templorum ab infima plebe completi erant. a qua plausu maximo cum esset mihi gratulatio significata, similis et frequentia et plausus me usque ad Capitolium celebravit in foroque et in ipso Capitolio miranda multitudo fuit.

\ciceroSection{6} postridie in senatu qui fuit dies Nonarum Septembr. senatui gratias egimus. eo biduo cum esset annonae summa caritas et homines ad theatrum primo, deinde ad senatum concurrissent, impulsu Clodi mea opera frumenti inopiam esse clamarent, cum per eos dies senatus de annona haberetur et ad eius procurationem sermone non solum plebis verum etiam bonorum Pompeius vocaretur idque ipse cuperet multitudoque a me nominatim ut id decernerem postularet, feci et accurate sententiam dixi. cum abessent consulares, quod tuto se negarent posse sententiam dicere, praeter Messallam et Afranium, factum est senatus consultum in meam sententiam, ut cum Pompeio ageretur ut eam rem susciperet lexque ferretur. quo senatus consulto recitato cum populus more hoc insulso et novo plausum meo nomine recitando dedisset, habui contionem. omnes magistratus praesentes praeter unum praetorem et duos tribunos pl. dederunt. postridie senatus frequens et omnes consulares nihil Pompeio postulanti negarunt. ille legatos quindecim cum postularet, me principem nominavit et ad omnia me alterum se fore dixit. legem consules conscripserunt qua Pompeio per quinquennium omnis potestas rei frumentariae toto orbe terrarum daretur, alteram Messius qui omnis pecuniae dat potestatem et adiungit classem et exercitum et maius imperium in provinciis quam sit eorum qui eas obtineant. illa nostra lex consularis nunc modesta videtur, haec Messi non ferenda. Pompeius illam velle se dicit, familiares hanc. consulares duce Favonio fremunt; nos tacemus et eo magis quod de domo nostra nihil adhuc pontifices responderunt. qui si sustulerint religionem, aream praeclaram habebimus; superficiem consules ex senatus consulto aestimabunt; sin aliter, †demolientur, suo† nomine locabunt, rem totam aestimabunt.

\ciceroSection{8} ita sunt res nostrae, ut ín secundis flúxae, ut in advorsís bonae. in re familiari valde sumus, ut scis, perturbati. praeterea sunt quaedam domestica quae litteris non committo. quin tum fratrem insigni pietate, virtute, fide praeditum sic amo ut debeo. te exspecto et oro ut matures venire eoque animo venias ut me tuo consilio egere non sinas. alterius vitae quoddam initium ordimur. iam quidam qui nos absentis defenderunt incipiunt praesentibus occulte irasci, aperte invidere. vehementer te requirimus.
